\documentclass[10pt]{article}
\usepackage[margin=20mm]{geometry}
\usepackage{amsmath,amssymb,mathtools,fontspec,unicode-math,xurl,hyperref}
\IfFontExistsTF{DejaVu Serif}{\setmainfont{DejaVu Serif}}{\setmainfont{Latin Modern Roman}}
\setmathfont{latinmodern-math.otf}
\setlength{\parindent}{0pt}\setlength{\parskip}{6pt}\sloppy
\title{Technical Review of Perfectoid Signature
and Perfectoid Hilbert–Kunz Multiplicity}
\author{AI-assisted technical review}
\date{}
\begin{document}
\maketitle
Abstract: This initial review concerns the 88-page author manuscript arXiv:2209.04046v4. It examines the definitions, main theorem dependencies, regularity and positivity criteria, finite-map transformation rule, and applications. One major proof gap is established in the global geometric application: the specified \(p\)-adically completed graded absolute integral closure cannot receive the asserted map from the completed cone. A separate minor calculation error affects two normalized-length constants; their correct values still supply the required positive bounds. Neither finding is a counterexample to a main theorem. The local fundamental-group argument is not shown to fail. External foundational results are used as cited dependencies, with the limits of this review stated below. Scientific adjudication is pending.\par
\section*{Paper identity and inspected scope}
The reviewed paper is Perfectoid signature, perfectoid Hilbert–Kunz multiplicity, and an application to local fundamental groups, by Hanlin Cai, Seungsu Lee, Linquan Ma, Karl Schwede, and Kevin Tucker [paper]. Its journal identifier is DOI `10.1007/s00222-025-01399-w`. The inspected object is specifically arXiv version 4, dated 3 July 2025, not an asserted copy of the journal version of record. All page numbers below are physical PDF pages, which agree with printed pages in this 88-page manuscript.\par
The bounded review covered the mathematical text of Sections 1–8 through the source files and PDF text, with deeper reconstructions of the normalized-length arguments, regularity characterization, BCM positivity, rational signature, finite-map transformation, and Section 7 applications. Page-image inspections are separately recorded. This is not a claim that every external cited theorem or every sentence has been independently proved. The derived and prismatic foundations, existence and weak functoriality of perfectoid big Cohen–Macaulay algebras, and several cited local-cohomology and cover-comparison results remain accepted dependencies.\par
No separate supplement or original computational implementation was obtained in the bounded acquisition search. This is a theoretical paper; absence of code is not a mathematical finding. The acquisition manifests identify the exact files consulted. A bounded primary-source search did not identify an author correction for the two findings; this does not establish novelty or absence of a correction elsewhere.\par
\section*{Definitions, quantifiers, and principal checks}
Let \((R,\mathfrak m,k)\) be a complete Noetherian local domain of mixed characteristic \((0,p)\), dimension \(d\), with perfect residue field. The paper fixes a finite coefficient-ring inclusion \[A=W(k)\mathopen{[\![}x_2,\ldots,x_d\mathclose{]\!]}\subset R,
 \qquad
 A_n=A[p^{1/p^n},x_2^{1/p^n},\ldots,x_d^{1/p^n}],\] and its \(p\)-adically completed tower \(A_\infty\). Normalized length on a finite-level module is scaled by \(p^{-nd}\). The relevant algebra is \(B=(R\otimes_A A_\infty)_{\mathrm{perfd}}\), not its image in a chosen absolute integral closure. For an \(\mathfrak m\)-primary ideal \(J\) and the non-splitting ideal \(I_\infty\subset B\), the invariants are \[e_{\mathrm{perfd}}^{\underline{x}}(J)=\lambda_\infty(B/JB),
 \qquad
 s_{\mathrm{perfd}}^{\underline{x}}(R)=\lambda_\infty(B/I_\infty).\] Their dependence on the chosen tower is retained throughout. Parameter independence is posed as Conjecture 8.1, not established by this paper. In particular, the class-group bound in Theorem 7.0.12 uses some choices of parameters; it must not be read as proving independence for arbitrary choices.\par
The regularity checks reconstructed the chain from the rank/multiplicity comparison and parameter limit closure to \(e_{\mathrm{perfd}}^{\underline{x}}(R)\geq1\), the equality criterion, and the implication \(s_{\mathrm{perfd}}^{\underline{x}}(R)=1\Rightarrow R\) regular. The final regularity argument relies on the cited homological criterion; its external proof was not re-audited. The distinction between \(B\) and its image is essential even for ramified regular rings and is correctly emphasized in the source.\par
For positive signature, the converse implication is quantified as \[\forall\,0\ne g\in R,\quad \exists e,\quad
 \forall\,\text{perfectoid BCM }R^+\text{-algebras }C:
 \quad R\xrightarrow{\,1\mapsto g^{1/p^e}\,}C
 \text{ is pure}.\] The exponent can depend on \(g\); it is not one exponent for all \(g\). The domination lemma and essential socle of the injective hull support uniformity over \(C\). The forward implication from BCM-regular pairs retains the effective boundary and log-\(\mathbb Q\)-Cartier assumptions. For rational signature, the stabilized tail-span argument in Theorem 5.4.10 supplies a positive valuation threshold uniform over all socle injections. Pointwise positivity alone would not justify the infimum. Problem 2 changes its numerical lower bound, not that uniformity.\par
The finite-map check retained normality, splitting, quasi-étaleness, the compatible tower, and the residue-field factor in Theorem 6.0.1. Trace adjunction and the two-sided normalized-length estimates give the stated transformation under these assumptions. Lemma 7.0.1 supplies splitting and ascent of weak BCM-regularity; together with Theorem 6.0.1 it supports the local cover-degree and fundamental-group bound in Theorem 7.0.2. These local arguments are not shown to be affected by the gap below. No stronger result without those hypotheses is inferred.\par
\section*{Substantive findings}
Problem 1 (The completed cone does not map to the specified target).  The natural factorization used in Claim 7.0.5, p. 75, is impossible for the algebra \(T\) defined there. This is a major proof gap in the derivation of the global geometric application, Theorem 7.0.4.\par
\subsection*{Claim, hypotheses, and dependency}
Theorem 7.0.4 concerns a normal integral scheme \(Y\), projective over \(R=H^0(Y,\mathcal O_Y)\), with equidimensional fibers. Here \(R\) is a complete Noetherian local domain with algebraically closed residue field. An effective \(\mathbb Q\)-divisor \(\Delta\) satisfies \(-(K_Y+\Delta)\) ample, and for every effective Cartier divisor \(D>0\) there is \(\epsilon>0\) such that \((Y,\Delta+\epsilon D)\) is globally \(+\)-regular. The conclusion is finiteness of the étale fundamental group of every open subset with complement of codimension at least two.\par
Equation reference: eq:badmap. In the proof, the chosen section ring \(S'\) is finite over \[A''=A'[x_0,\ldots,x_r],\qquad
 A'=W(k)\mathopen{[\![}y_1,\ldots,y_s\mathclose{]\!]}.\] Write \(\mathfrak n\) for its homogeneous maximal ideal and \(S=\widehat{S'}_{\mathfrak n}\). Claim 7.0.5 asserts BCM-regularity of \((S,\Delta_S)\) and hence positive perfectoid signature. After another graded finite extension and completion \(S_1\), the proof asserts a natural factorization \[
 S\longrightarrow S_1\longrightarrow
 T:=\widehat{S'^{+,\mathrm{GR}}}^{\,p}.\] It explicitly defines the last completion as \(p\)-adic. This factorization supplies the tensor product and local-cohomology map used to invoke Proposition 5.2.3. The subsequent reduction of covers of \(Y\) to covers of \(S\) needs the resulting positive signature.\par
\subsection*{Exact obstruction}
Equation reference: eq:residue. Put \(C=S'^{+,\mathrm{GR}}\). Under the stated definition, \(C\) is an integral extension of \(S'\), so it is integral over \(A''\). Consider the maximal ideal \[\mathfrak a=(p,y_1,\ldots,y_s,x_0-1,x_1,\ldots,x_r)
 \subset A'' .\] It is proper: its quotient is \(k\). Lying over gives a prime \(\mathfrak q\subset C\) contracting to \(\mathfrak a\). Consequently \(p\in\mathfrak q\) and \(1-x_0\in\mathfrak q\). There is a unital map to a nonzero ring \[
 T=\varprojlim_m C/p^mC\longrightarrow C/pC
 \longrightarrow C/\mathfrak q .\] Thus the image of \(1-x_0\) in \(T\) is not a unit.\par
By contrast, \(x_0\in\mathfrak n\), so \(1-x_0\) is a unit in the complete local ring \(S\); its inverse is the convergent \(\mathfrak n\)-adic series \(\sum_{j\geq0}x_0^j\). Every unital \(S'\)-algebra map \(S\to T\) would preserve this unit and its inverse, contradicting [eq:residue] (\#eq:residue). Therefore the natural \(S'\)-algebra map required by [eq:badmap] (\#eq:badmap) does not exist. The argument uses precisely the finite integral extension and completion specified in the proof. It does not assume failure of global \(+\)-regularity or require a counterexample to the geometric theorem.\par
\subsection*{Alternative readings and counterevidence}
The PDF image on p. 75 explicitly distinguishes completion at the homogeneous maximal ideal from the \(p\)-adic completion defining \(T\). The source TeX agrees; this is not an OCR substitution. The same paper does not define \(S'^{+,\mathrm{GR}}\) here to include prior localization at \(\mathfrak n\). A localization or a suitable completion after base change could remove this particular unit obstruction, but would change the target and require justification of the properties used next. Merely assuming that a formal power series converges \(p\)-adically does not fix the problem: \(x_0\) is not \(p\)-adically topologically nilpotent, as the quotient in [eq:residue] (\#eq:residue) also shows.\par
The cone construction and the local fundamental-group theorem may still yield the intended result with an appropriate auxiliary algebra. This possibility is counterevidence against concluding that Theorem 7.0.4 is false; it does not validate its displayed factorization.\par
\subsection*{Impact and repair status}
The failure is in a necessary map for Claim 7.0.5, so the stated local-cohomology injection cannot be used with this \(T\). The affected chain is Claim 7.0.5 \(\Rightarrow\) the cone positivity step \(\Rightarrow\) Theorem 7.0.4, including the proof-based relaxation in Remark 7.0.6. This is a material gap in an advertised geometric application, warranting Major at the claim level. It is not asserted to damage Theorems A–G, the local fundamental-group bound, or the torsion class-group argument.\par
A repair must construct a target receiving \(S_1\), establish the required perfectoid structure and compatible roots, and prove that the splitting/local-cohomology information from the graded covers survives the intervening operations. No complete replacement argument has been verified in this review. Author adoption is unknown.\par
Problem 2 (Total-degree volume is missing from normalized length).  The proof of Lemma 4.0.13, p. 32, uses an incorrect monomial-count lower bound. The final display of Theorem 5.4.11, p. 60, similarly assigns an incorrect scalar to a valuation-ideal quotient. This is a minor local calculation error, with a verified correction preserving both qualitative conclusions.\par
\subsection*{Claim and assumptions}
Definition 4.0.11, p. 31, extends the \(\mathfrak m_A\)-adic order to the tower by total monomial degree, with \(v(p)=v(x_i)=1\). For \(\epsilon>0\), the lemma sets \(I_n=\{a\in A_n:v(a)\geq\epsilon\}\) and uses \[\operatorname{length}_A(A_n/I_n)\geq p^{nd}\epsilon^d.\] At the end of Theorem 5.4.11, the paper instead calculates \(\lambda_\infty(A_\infty/\{v\geq p^{-e_0}\})=p^{-e_0}\). Both quantities must respect the dimension-\(d\) normalization.\par
\subsection*{Independent exact calculation}
Equation reference: eq:binomial, eq:volume. Let \(q=p^n\) and \(\mathfrak m_n=(p^{1/q},x_2^{1/q},\ldots,x_d^{1/q})\). The ring \(A_n\) is regular local with residue field \(k\), and \(v|_{A_n}=q^{-1}\operatorname{ord}_{\mathfrak m_n}\). For \(a>0\), put \(I_a=\{b\in A_\infty:v(b)\geq a\}\). Then \[
 I_a\cap A_n=\mathfrak m_n^{\lceil qa\rceil},
 \qquad
 \operatorname{length}_A\bigl(A_n/(I_a\cap A_n)\bigr)
   =\binom{\lceil qa\rceil+d-1}{d}.\] Indeed, degree-\(j\) monomials in \(d\) regular parameters number \(\binom{j+d-1}{d-1}\), and summing for \(0\leq j<\lceil qa\rceil\) gives [eq:binomial] (\#eq:binomial). Lengths over \(A\) and \(A_n\) agree because the residue fields agree. The finite-level formula for cyclic-module normalized length gives \[
 \lambda_\infty(A_\infty/I_a)
 =\lim_{n\to\infty}p^{-nd}
     \binom{\lceil p^n a\rceil+d-1}{d}
 =\frac{a^d}{d!}.\] For any \(a>0\), choose an integer \(N\geq a\). Then \(p^N A_\infty\subseteq I_a\). The uncompleted tower and its completion agree modulo \(p^N\), so completion does not change the quotient or this calculation.\par
For an explicit admissible instance, take \(R=A=\mathbb Z_2\mathopen{[\![}
x\mathclose{]\!]}\), \(d=2\), \(n=1\), and \(\epsilon=1\). Then \(I_1=\mathfrak m_1^2\) and the associated graded basis of \(A_1/I_1\) is \(1,2^{1/2},x^{1/2}\). Its length is \(3\), whereas the printed lower bound is \(4\). For the second formula choose \(d=2,p=2,e_0=1\). The actual limit is \(1/8\), whereas the printed value is \(1/2\): \[2^{-2n}\binom{2^{n-1}+1}{2}
   =\frac18+2^{-n-2}\longrightarrow\frac18 .\] The retained finite diagnostic independently enumerates monomials in 48 cases; all agree with [eq:binomial] (\#eq:binomial). That finite check supports, and does not replace, the general derivation above.\par
\subsection*{Counterevidence, correction, and bounded impact}
In dimension one the printed scalars agree with the correct formula. For \(d>1\), neither a change from floor to ceiling nor a different reading of the exponent removes the simplex-volume factor. The PDF images and the TeX agree on the disputed expressions. The valuation is explicitly total degree, not a coordinatewise minimum or a rescaled volume.\par
For Lemma 4.0.13, the correct lower bound \(p^{nd}\epsilon^d/d!\) remains strictly positive after normalization and yields the same contradiction to zero normalized length. In Theorem 5.4.11 the replacement is \[s_{\mathrm{rat},\mathrm{perfd}}^{\underline{x}}(R)
   \geq \frac{p^{-de_0}}{d!}>0.\] Together with the type bound, the relative signature is at least \(p^{-de_0}/(t\,d!)\) when \(t\) is the Cohen–Macaulay type. The exponent \(e_0\) is uniform over socle injections by Theorem 5.4.10, so taking the infimum remains valid.\par
The erroneous numerical values are auxiliary constants in proofs; the main theorem statements assert qualitative vanishing/positivity equivalences, not those numerical lower bounds. The corrected calculation restores the exact needed implications under unchanged hypotheses. Minor is assigned on that demonstrated dependency and bounded impact, not merely because the repair is short. No author adoption is documented.\par
\section*{Overall assessment and conclusion}
The highest supported assessment is Major, restricted to the invalid target-map construction in the global geometric application. The count issue is separately Minor. Neither is a proved false theorem, and the review does not establish wholesale invalidity of the paper. The inspected local fundamental-group and transformation arguments retain their stated assumptions. Unexamined external dependencies limit any positive conclusion about the whole manuscript.\par
Item | Assessment | Repair verification | Author adoption\par
Completed-cone target map | Major | No complete repair verified | Unknown\par
Normalized-length count | Minor | Corrected local implications checked | Unknown\par
Other inspected chains | No further issue established | Scoped checks only | Not applicable\par
Scientific adjudication remains pending, with no human reviews recorded. Initial delivery is not independent scientific acceptance.\par
9 H. Cai, S. Lee, L. Ma, K. Schwede, and K. Tucker, Perfectoid signature, perfectoid Hilbert–Kunz multiplicity, and an application to local fundamental groups, arXiv:2209.04046v4, 3 July 2025, 88 pp. https://arxiv.org/abs/2209.04046v4. Exact acquired PDF and author TeX inspected. Journal DOI: https://doi.org/10.1007/s00222-025-01399-w.\par
\section*{Supporting checks}
The bound main PDF has SHA256 `0001a9a56839319b5d36c79fcd0dedc9dd66ac0b38029773d923fa4037f944fa`. The source manifest has SHA256 `2efc2d0bdb277b11afde377828066e65a290ae53d97a208b4fa32cc441faf193`.\par
The monomial diagnostic uses exact integer arithmetic with \(d=1,2,3,4\), \(p=2\), \(n=1,2,3,4\), and \(a=1/2,1,3/2\). It enumerates tuples of nonnegative exponents with sum less than \(\lceil p^n a\rceil\), compares their count with the binomial formula, and separately evaluates the printed lower bound. All 48 exact-count comparisons pass; 33 violate the printed bound. The general result is established by [eq:binomial] (\#eq:binomial)–[eq:volume] (\#eq:volume), not by extrapolation from those diagnostics.\par
The current conversation performed the mathematical synthesis and report writing. The controller-reported UI label and the runtime’s self-description are recorded separately in provenance; no equivalence between them is asserted. No author code was run, no author source was compiled, and no external dependency was silently substituted. No author has approved or adopted the proposed correction in the available record.\par
\end{document}
